% !TeX root = ../main.tex

\chapter{Basics of Quantum State Discrimination}
\label{chap:basics_qsd}

% Revised by Gemini 2.5 Flash
This chapter reviews fundamental concepts in quantum state discrimination.

Throughout this thesis, we adopt the standard convention in quantum
communication of an information sender (Alice) and a receiver (Bob).
For a foundational understanding of these concepts, including scenarios
involving eavesdroppers, readers are referred to Sections 2.2.4 and 12.1 of
\cite{qcqi} or Chapter 3 of \cite{10.5555/3240076}.

The core problem of quantum state discrimination involves Alice preparing a
quantum state from a predefined set of states known to both herself and Bob.
Bob's objective is then to perform a measurement to identify which specific
state Alice transmitted \cite{qcqi}. A critical challenge arises with
non-orthogonal states: these states cannot be perfectly discriminated.
This means no measurement exists that can reliably identify each state in a
non-orthogonal set with unit probability \cite{qcqi}.

\section{The quantum state discrimination problem}
% Revised by Gemini 2.5 Flash
Let $S_{\ket{\psi}} \equiv \{\ket{\psi_{1}}, \ket{\psi_{2}}, \dots, \ket{\psi_{k}}\}$
denote the predefined set of quantum states shared between Alice and Bob.
For quantum state discrimination, the states in $S_{\ket{\psi}}$ are typically 
non-orthogonal, meaning their inner product is non-zero for distinct states:
\begin{align}
    \braket{\psi_{i}}{\psi_{j}} \neq 0, \forall \; i \neq j, 1 \leq i, j \leq k.
\end{align}
Let $S_{prior} \equiv \{p_{1}, p_{2}, \dots, p_{k}\}$ be the set of prior 
probabilities associated with each state in $S_{\ket{\psi}}$. The prior 
probability $p_{i}$ represents the likelihood that Alice prepares and 
transmits the state $\ket{\psi_{i}}$ to Bob.

Given the predefined set of states $S_{\ket{\psi}}$ and their prior 
probabilities $S_{prior}$, the quantum state discrimination problem focuses on 
determining the optimal measurement strategy for Bob to identify the quantum 
state prepared by Alice. Numerous figures of merit have been proposed to 
quantify performance in quantum state discrimination \cite{RN65}.
The most common figure of merit is the minimization of measurement error, 
which is equivalent to the maximization of the success probability of 
discrimination.

To illustrate these concepts and different discrimination methods, consider the following example:

Let $k = 2$ with states $\ket{\psi_{1}} = \ket{0}$ and $\ket{\psi_{2}} = \dfrac{\ket{0}+\ket{1}}{\sqrt{2}}$,
and prior probabilities $S_{prior} = \{p_{1}, p_{2}\} = \{0.5, 0.5\}$.
These two states are non-orthogonal, as demonstrated by their inner product:
$\braket{\psi_{1}}{\psi_{2}} = \dfrac{1}{\sqrt{2}} \neq 0$.

\section{Projective measurement}
% Revised by Gemini 2.5 Flash
A projective measurement is defined by a set of $k$ measurement operators, $S_{\Pi_{proj}} \equiv \{\Pi_{1}, \Pi_{2}, \dots, \Pi_{k}\}$, which satisfy the following properties:
\begin{enumerate}
    \item \textbf{Completeness Relation:} $\sum_{i = 1}^{k}\Pi_{i} = I$ (the identity operator).
    \item \textbf{Orthogonality:} $\Pi_{i}\Pi_{j} = \delta_{ij}\Pi_i$, where $\delta_{ij}$ is the Kronecker delta.
    \item \textbf{Hermiticity and Projectors:} Each $\Pi_i$ is a Hermitian operator ($\Pi_i^\dagger = \Pi_i$) and satisfies $\Pi_i^2 = \Pi_i$, meaning they are projectors onto orthogonal subspaces.
\end{enumerate}

For the example states introduced previously, a standard projective measurement in the computational basis is given by:
\begin{align}
    &\Pi_{1} = \ket{0}\bra{0} = \begin{bmatrix} 1 \\ 0 \end{bmatrix} \begin{bmatrix} 1 & 0 \end{bmatrix} = \begin{bmatrix} 1 & 0 \\ 0 & 0 \end{bmatrix} \\
    &\Pi_{2} = \ket{1}\bra{1} = \begin{bmatrix} 0 \\ 1 \end{bmatrix} \begin{bmatrix} 0 & 1 \end{bmatrix} = \begin{bmatrix} 0 & 0 \\ 0 & 1 \end{bmatrix}
\end{align}

The probability of success ($P_{succ}$) for this measurement can be calculated as follows:
\begin{align}
    &\bra{\psi_{1}} \Pi_{1} \ket{\psi_{1}} = 1, \quad \bra{\psi_{1}} \Pi_{2} \ket{\psi_{1}} = 0 \\
    &\bra{\psi_{2}} \Pi_{1} \ket{\psi_{2}} = 0.5, \quad \bra{\psi_{2}} \Pi_{2} \ket{\psi_{2}} = 0.5 \\
    &P_{succ} = p_{1} \bra{\psi_{1}} \Pi_{1} \ket{\psi_{1}} + p_{2} \bra{\psi_{2}} \Pi_{2} \ket{\psi_{2}} \\
    &= 0.5 \times 1 + 0.5 \times 0.5 = 0.75
\end{align}

A measurement error occurs when Bob's measurement outcome does not correspond to the state Alice actually sent.
In this example, an error arises when Alice sends $\ket{\psi_{2}}$ but Bob's measurement yields outcome $\Pi_{1}$ (i.e., $\ket{0}$).
This specific projective measurement strategy thus only incurs error when
discriminating $\ket{\psi_{2}}$, as $\bra{\psi_{2}} \Pi_{1} \ket{\psi_{2}} = 0.5$,
meaning there's a 50\% chance of incorrectly inferring $\ket{\psi_1}$ when $\ket{\psi_2}$ was sent.

\section{Minimum Error State Discrimination}
Instead of measuring with a trivial projective measurement,
the first proposed method is minimum error state discrimination (MESD),
also known as minimum error discrimination (MED).
The principle of MED is to find a POVM of $l$ elements that minimizes the
overall measurement error
\begin{align}
     & P_{err} = \Sigma_{i = 1, i \neq j}^{l} p_{i} \bra{\psi_{i}} \Pi_{j} \ket{\psi_{i}} \\
     & = 1 - \Sigma_{i = 1}^{l} p_{i} \bra{\psi_{i}} \Pi_{i} \ket{\psi_{i}}               \\
     & = 1 - P_{succ}
\end{align}
$l$ here is exactly the number of states in $S_{\ket{\psi}}$.
The POVM is denoted by $S_{\Pi} \equiv \{\Pi_{1}, \Pi_{2}, \dots, \Pi_{l}\}$.

For $l = 2$, the optimal probability of success $P_{succ}$ can be analytically
solved, which is also named the Helstrom measurement \cite{Barnett:09}.
\begin{align}
     & \Sigma_{i=1}^{l} \Pi_{i} = \Pi_{1} + \Pi_{2} = I                          \\
     & P_{succ} = \Sigma_{i = 1}^{l} p_{i} \bra{\psi_{i}} \Pi_{i} \ket{\psi_{i}} \\
     & = p_{1} \bra{\psi_{1}} \Pi_{1} \ket{\psi_{1}}
    + p_{2} \bra{\psi_{2}} \Pi_{2} \ket{\psi_{2}}                                \\
     & = p_{1} \bra{\psi_{1}} \Pi_{1} \ket{\psi_{1}}
    + p_{2} \bra{\psi_{2}} (I - \Pi_{1}) \ket{\psi_{2}}                          \\
     & = p_{2} + \operatorname{Tr} \left[
        \left( p_{1} \ket{\psi_{1}}\bra{\psi_{1}}
        - p_{2} \ket{\psi_{2}}\bra{\psi_{2}} \right) \Pi_{1}
        \right] \label{eq:psucc}
\end{align}

From Equation \ref{eq:psucc}, the maximal probability happens when $\Pi_{1}$
projects to the positive eigenspace of
$p_{1} \ket{\psi_{1}}\bra{\psi_{1}} - p_{2} \ket{\psi_{2}}\bra{\psi_{2}}$.

For our example,
\begin{align}
     & p_{1} \ket{\psi_{1}}\bra{\psi_{1}} - p_{2} \ket{\psi_{2}}\bra{\psi_{2}}
    = 0.5 \times \begin{bmatrix}
                     1 & 0 \\ 0 & 0
                 \end{bmatrix}
    - 0.5 \times \begin{bmatrix}
                     0.5 & 0.5 \\ 0.5 & 0.5
                 \end{bmatrix}                                         \\
     & = 0.5 \times \begin{bmatrix}
                        0.5 & -0.5 \\ -0.5 & -0.5
                    \end{bmatrix}
    \equiv 0.5 A
\end{align}

The eigenvectors $v_{1}, v_{2}$ and eigenvalues $\lambda_{1}, \lambda_{2}$ of $A$ are
\begin{align}
     & v_{1} = \begin{bmatrix} -1 - \sqrt{2} & 1   \end{bmatrix} ^{T},
    \lambda_{1} =\frac{\sqrt{2}}{2}                                     \\
     & v_{2} = \begin{bmatrix}  -1 + \sqrt{2} & 1   \end{bmatrix} ^{T},
    \lambda_{2} =\frac{-\sqrt{2}}{2}
\end{align}

We choose $v_{1}$ since we want the one with positive eigenvalue.
The resulting measurement operators are
\begin{align}
     & \Pi_{1} = \frac{v_{1}v_{1}^{T}}{(-1 - \sqrt{2})^{2} + 1^{2}}
    = \begin{bmatrix}
          3 + 2\sqrt{2} & -1 - \sqrt{2} \\
          -1 - \sqrt{2} & 1
      \end{bmatrix}
    / (4 + 2\sqrt{2})                                               \\
     & = (2 - \sqrt{2})\begin{bmatrix}
                           3 + 2\sqrt{2} & -1 - \sqrt{2} \\
                           -1 - \sqrt{2} & 1
                       \end{bmatrix}
    / 4                                                             \\
     & = \begin{bmatrix}
             (2 + \sqrt{2})/4 & - \sqrt{2}/4     \\
             - \sqrt{2}/4     & (2 - \sqrt{2})/4
         \end{bmatrix}                        \\
     & \Pi_{2} = I - \Pi_{1} = \begin{bmatrix}
                                   (2 - \sqrt{2})/4 & \sqrt{2}/4       \\
                                   \sqrt{2}/4       & (2 + \sqrt{2})/4
                               \end{bmatrix}
\end{align}

The probability of success for MED is
\begin{align}
     & P_{succ} = p_{1} \bra{\psi_{1}} \Pi_{1} \ket{\psi_{1}}
    + p_{2} \bra{\psi_{2}} \Pi_{2} \ket{\psi_{2}}                  \\
     & = 0.5 \times (2 + \sqrt{2})/4 + 0.5 \times (2 + \sqrt{2})/4 \\
     & = (2 + \sqrt{2})/4 = 0.85355\dots > 0.75
\end{align}

\section{Unambiguous Quantum State Discrimination}

Unambiguous quantum state discrimination (UQSD, USD) \cite{RN493, RN494, RN495}
provided a different view, where the goal is to find a POVM of $l+1$
elements, such that the measurement may obtain inconclusive results,
while the conclusive results will not be erroreous at all.

For example,


The B92 protocol used USD and is one of the simplest quantum key distribution
protocols. In the protocol, Bob uses USD to identify Alice's messages.
Since it is easy to implement, B92 is often used in experiments instead of
BB84.
However, since Eve can also perform USD, she can resend Bob some states
without being detected, which is also called the zero-error attack. Even
though there are some countermeasures for USD attacks \cite{RN226}, BB84 with
decoy states is still the mainstream.

previous works
Theoretical bounds
Physics experiments

\section{Optimal POVM for UQSD}
\label{sec:optimal_povm_for_uqsd}

\cite{RN23}
In this section, we will review the work that defined the optimal
POVM for UQSD.
The work is limited to linearly independent pure states and

Its problem is that some probabilities can be zero or near-zero, causing the
discrimination scheme to throw away some states.

We mention this work because it is fast when everything just works.

\subsection{Entangled Quantum States with Orthogonal Components}
(Over-simplified example)
We define a trio of two-qubit entangled states, each expressed as
a linear combination of basis states in the computational basis.
The states are given by:
\begin{align}
    \ket{\psi_1} & = \frac{1}{\sqrt{1 + a_{1}^{2}}} \left( \ket{00} + a_{1} \ket{11} \right) , \\
    \ket{\psi_2} & = \frac{1}{\sqrt{1 + a_{2}^{2}}} \left( \ket{01} + a_{2} \ket{11} \right) , \\
    \ket{\psi_3} & = \frac{1}{\sqrt{1 + a_{3}^{2}}} \left( \ket{10} + a_{3} \ket{11} \right)
\end{align}
where $a_{i} \in \mathbb{R}, 0 < a_{i} < 1$ for $i \in \{1, 2, 3\}$.
This construction ensures that each state is a valid pure quantum state,
and the orthogonal components ($\ket{00}, \ket{01}, \ket{10}$) weigh more in the linear combination.
It is lucid that the state vectors are non-orthogonal, which makes perfect discrimination impossible,
and they are also linearly independent, which is a necessary assumption for UQSD to exist.

For example, given $a_{1} = 0.2, a_{2} = 0.5$, and $a_{3} = 0.7$.
We can express the states in state vectors, where the floating point numbers are rounded to 4 digits.
\begin{align*}
    \ket{\psi_1} & = \begin{pmatrix}
                         0.9806+0.j &  & 0.    +0.j &  & 0.    +0.j &  & 0.1961+0.j
                     \end{pmatrix}^{T} \\
    \ket{\psi_2} & = \begin{pmatrix}
                         0.    +0.j &  & 0.8944+0.j &  & 0.    +0.j &  & 0.4472+0.j
                     \end{pmatrix}^{T} \\
    \ket{\psi_3} & = \begin{pmatrix}
                         0.    +0.j &  & 0.    +0.j &  & 0.8192+0.j &  & 0.5735+0.j
                     \end{pmatrix}^{T}
\end{align*}

In the following we will show the process of deriving the semidefinite
programming problem for the optimal POVM for pure quantum states.

% Morning stops here
The pseudo-inverse is
\begin{align*}
    Q_{1}
\end{align*}

The POVM is
\begin{align*}
    \Pi_{1} = \begin{bmatrix}
                  1
              \end{bmatrix} \\
    \Pi_{2} = \begin{bmatrix}
                  1
              \end{bmatrix} \\
    \Pi_{3} = \begin{bmatrix}
                  1
              \end{bmatrix} \\
    \Pi_{?} =
\end{align*}

% Noise model

\begin{itemize}
    \item POVM results with annotation
    \item solve time?
    \item
\end{itemize}

% \subsubsection{$a=b=c$}
% \subsubsection{$a>b>c$}
\subsubsection{Parameter search}
According to our SDP formulation, we can trade off
error by adjusting the parameters. In this subsection we
will discuss the results of different parameter search schemes.




My research framework
Experiment setup
1. Qiskit

2. Discriminate two-to-four-qubit state

3. circuit

According to previous publications, it is possible to find the optimal probability
of the unambiguous quantum state discrimination by a semidefinite programming approach [].

We apply a similar approach and convert it into a quantum circuit.
A convex optimization problem is formed by the target states,
which is a set of linearly independent pure quantum states.

By optimal, it means the measurement gives the smallest possible
probability of obtaining an inclusive result for unambiguous discrimination
\cite{RN23}


% Generated by Google Gemini 2.5 Flash
\section{The Holevo Bound}

The Holevo bound is a foundational result in quantum information theory that
establishes a fundamental limit on the amount of classical information that
can be extracted from an ensemble of quantum states.
It provides an upper bound on the classical mutual information between the
preparation of quantum states and the outcomes of a measurement performed on them.

Consider an ensemble of quantum states $\mathcal{E} = \{p_x, \rho_x\}$, where
the state $\rho_x$ is prepared with probability $p_x$.
Let the measurement be described by a set of positive-operator valued measure
(POVM) elements $\{M_y\}$, which results in a measurement outcome $y$.
The classical mutual information $I(X;Y)$ between the choice of state $X$ and
the measurement outcome $Y$ is given by
\begin{equation}
    I(X; Y) = H(Y) - H(Y|X)
\end{equation}
where $H(Y)$ and $H(Y|X)$ are the classical Shannon entropy and conditional entropy, respectively.

The Holevo bound states that this classical mutual information is limited by
the Holevo information $\chi(\mathcal{E})$ of the ensemble.
The inequality is given by:
\begin{equation}
    I(X; Y) \leq S(\bar{\rho}) - \sum_x p_x S(\rho_x) = \chi(\mathcal{E})
\end{equation}
where $S(\cdot)$ is the von Neumann entropy and $\bar{\rho} = \sum_x p_x \rho_x$
is the average density matrix of the ensemble.

This inequality demonstrates that while an ensemble of non-orthogonal quantum
states may appear to carry a large amount of information, a single measurement
is fundamentally limited in its ability to extract all of it.
The Holevo bound reveals a crucial distinction between classical and quantum
information, highlighting the inaccessibility of information encoded in non-orthogonal states.




\section{Scalability/Complexity}


state exclusion, hypothesis
